\chapter{H}
\label{chap:h}

\term{Haar Measure}
A translation-invariant measure on a locally compact topological group. Such a measure exists and is unique up to a positive scalar factor.

\term{Hadamard Factorization Theorem}
The theorem stating that any entire function of finite order can be expressed as a product involving its zeros times an exponential polynomial. It generalises the Weierstrass product theorem.

\term{Hadamard Matrix}
A square matrix whose entries are either $+1$ or $-1$ and whose rows are mutually orthogonal. Such matrices are used in error-correcting codes and signal processing.

\term{Hadamard Product}
The element-wise product of two matrices of the same dimensions, where $(A \circ B)_{ij} = A_{ij} B_{ij}$. It is also known as the Schur product.

\term{Hahn-Banach Theorem}
A fundamental result in functional analysis stating that a bounded linear functional on a subspace of a normed vector space can be extended to the whole space without increasing its norm.

\term{Hairy Ball Theorem}
Every continuous tangent vector field on an even-dimensional sphere must vanish somewhere; equivalently, there is no continuous non-vanishing field of tangent vectors on $S^{2n}$. It is often summarized as ``you cannot comb a hairy ball flat.'' A consequence is that at any moment the wind must be calm at some point on Earth. Odd-dimensional spheres do admit non-vanishing tangent vector fields.

\term{Half-plane}
A region in $\mathbb{R}^2$ bounded by a line, consisting of all points on one side of the line.

\term{Half-space}
A region in $\mathbb{R}^n$ bounded by a hyperplane.

\term{Hall's Marriage Theorem}
A combinatorial theorem giving a necessary and sufficient condition for a finite family of sets to have a system of distinct representatives: the union of any $k$ sets must contain at least $k$ elements.

\term{Halting Problem}
The problem of deciding, given a program and its input, whether the program eventually halts. Turing proved that it is undecidable: no algorithm solves it in all cases. The proof uses a diagonalization argument, constructing a program that halts exactly when the purported decider claims it does not. It is the classic undecidable problem and the starting point for the theory of computability.

\term{Hamiltonian (Mechanics)}
A function $H(p, q, t)$ representing the total energy of a system. The evolution of the system is governed by Hamilton's equations: $\dot{q} = \frac{\partial H}{\partial p}$ and $\dot{p} = -\frac{\partial H}{\partial q}$.

\term{Hamiltonian Operator}
In quantum mechanics, the operator $\hat{H}$ corresponding to the total energy of the system. The time-independent Schrödinger equation is $\hat{H} \psi = E \psi$.

\term{Hamiltonian System}
A dynamical system governed by Hamilton's equations $\dot{q} = \frac{\partial H}{\partial p}$ and $\dot{p} = -\frac{\partial H}{\partial q}$.

\term{Hamming Code}
A family of linear error-correcting codes that can detect up to two-bit errors and correct one-bit errors. The (7,4) Hamming code is a prototypical example.

\term{Hamming Distance}
The number of positions at which the corresponding symbols of two strings of equal length are different. It is used in information theory to measure error rates.

\term{Hankel Transform}
An integral transform of the form $F(r) = \int_0^\infty f(t) J_n(rt)\, t\, dt$, where $J_n$ is a Bessel function of the first kind. It is the Fourier transform adapted to radial symmetry.

\term{Harmonic Analysis}
The study of how signals or functions can be represented as sums of basic waves (sinusoids), primarily via Fourier series and Fourier transforms.

\term{Harmonic Function}
A function $u$ that satisfies Laplace's equation $\Delta u = 0$. Such functions are smooth and satisfy the mean value property.

\term{Harmonic Mean}
The reciprocal of the arithmetic mean of the reciprocals of a set of values: $H = n / (\sum 1/x_i)$. It is typically used for average rates.

\term{Harmonic Number}
The sum of the reciprocals of the first $n$ positive integers: $H_n = \sum_{k=1}^n 1/k$. The harmonic numbers grow like $\ln n + \gamma$, where $\gamma$ is Euler's constant.

\term{Harmonic Series}
The divergent infinite series $\sum_{n=1}^\infty \frac{1}{n}$. Although the terms tend to zero, the sum grows logarithmically.

\term{Harnack's Inequality}
An inequality bounding the ratio of values of a positive harmonic function at two points in terms of a constant depending only on the domain. It implies that positive harmonic functions on connected domains satisfy a strong form of rigidity.

\term{Hausdorff Dimension}
A way of assigning a dimension to a set (including fractals) that can be non-integer. It is defined via the $s$-dimensional Hausdorff measure.

\term{Hausdorff Distance}
A measure of how far two subsets of a metric space are from each other. It is the maximum distance from a point in one set to the nearest point in the other set.

\term{Hausdorff Paradox}
A unit sphere $S^2$ can be decomposed into finitely many disjoint pieces which, by rotations alone, can be reassembled into two copies of the sphere (one possibly together with a countable set). It is a precursor of the Banach--Tarski paradox and similarly depends on the axiom of choice and non-measurable sets.

\term{Hausdorff Space}
A topological space in which any two distinct points can be separated by disjoint open neighborhoods. Most "natural" spaces in analysis and geometry are Hausdorff.

\term{Heat Equation}
The partial differential equation $\frac{\partial u}{\partial t} = \alpha \nabla^2 u$ describing the diffusion of heat through a medium.

\term{Heaviside Step Function}
A function $\text{H}(x)$ that is $0$ for $x < 0$ and $1$ for $x \geq 0$. It is used to model switches in circuit theory and is the integral of the Dirac delta function.

\term{Hedgehog Space}
A topological space formed by attaching several copies of the unit interval to a single common endpoint. It is used as a counterexample in the study of metric spaces.

\term{Heegner Number}
A number $d$ such that the imaginary quadratic field $\mathbb{Q}(\sqrt{-d})$ has class number 1. There are exactly nine such numbers: 1, 2, 3, 7, 11, 19, 43, 67, 163.

\term{Heisenberg Group}
A Lie group of $3 \times 3$ upper triangular matrices with ones on the diagonal. It is the simplest non-abelian nilpotent Lie group and is central to quantum mechanics.

\term{Helly's Theorem}
A result in convex geometry stating that if every $d+1$ elements of a collection of convex sets in $\mathbb{R}^d$ have a common intersection, then the entire collection has a common intersection.

\term{Hensel's Lemma}
A result in number theory stating that if a polynomial has a simple root modulo a prime $p$, then it has a unique root modulo $p^k$ for every $k$, obtained by successive lifting.

\term{Hereditary Property}
A property of a class of graphs or structures such that if a structure has the property, then all its substructures (e.g., induced subgraphs) also have it.

\term{Hereditary Ring}
A ring where every ideal is a projective module. In such rings, the global dimension is at most 1.

\term{Hermite Normal Form}
A canonical integer matrix form obtained by column operations, in which the matrix is upper triangular with non-negative off-diagonal entries. It is used in lattice theory and integer linear programming.

\term{Hermite Polynomials}
A sequence of orthogonal polynomials $H_n(x)$ that arise in the solution of the quantum harmonic oscillator and in probability theory.

\term{Hermitian Form}
A sesquilinear form $h$ on a complex vector space satisfying $h(u,v) = \overline{h(v,u)}$. A positive definite Hermitian form is an inner product.

\term{Hermitian Matrix}
A complex square matrix $A$ that is equal to its own conjugate transpose: $A = A^*$. All eigenvalues of a Hermitian matrix are real.

\term{Hessian Matrix}
The square matrix of second-order partial derivatives of a scalar-valued function. It is used to determine if a critical point is a local maximum, minimum, or saddle point.

\term{Heuristics}
Practical methods or "rules of thumb" used in problem solving that may not be guaranteed to find an optimal solution but often find a "good enough" one in a reasonable time.

\term{Hilbert Basis Theorem}
The theorem that the polynomial ring $R[x_1, \ldots, x_n]$ over a Noetherian ring $R$ is itself Noetherian.

\term{Hilbert Curve}
A space-filling curve that maps the unit interval $[0,1]$ onto the unit square $[0,1]^2$ continuously. It is a fractal curve.

\term{Hilbert Space}
A complete inner product space. They generalize Euclidean space to infinite dimensions and are the primary setting for quantum mechanics.

\term{Hilbert Transform}
An integral transform that shifts the phase of every frequency component of a signal by $90^\circ$. It is used to create the analytic signal in communications.

\term{Hilbert's Nullstellensatz}
A fundamental theorem in algebraic geometry relating ideals in polynomial rings to the zero-sets (varieties) they define. It states that every ideal in $k[x_1, \ldots, x_n]$ is the radical of the ideal of its zero-set.

\term{Hilbert's Tenth Problem}
The problem of determining whether a given Diophantine equation has a solution in integers. This was proven undecidable by Yuri Matiyasevich in 1970.

\term{Hilbert-Schmidt Norm}
The norm of a linear operator $T$ on a Hilbert space defined by $\|T\|_{HS}^2 = \sum \|Te_i\|^2$ for any orthonormal basis $\{e_i\}$.

\term{Hilbert-Schmidt Operator}
A bounded linear operator whose Hilbert-Schmidt norm is finite. Such operators are always compact.

\term{Hille-Yosida Theorem}
A theorem characterising the generators of strongly continuous semigroups of contractions on a Banach space in terms of resolvent bounds.

\term{H-index}
A metric that measures both the productivity and citation impact of a publication, where a scientist has index $h$ if $h$ of their papers have at least $h$ citations each.

\term{Hodge Star Operator}
A linear isomorphism between $k$-forms and $(n-k)$-forms on an $n$-dimensional oriented Riemannian manifold. It is used to define the adjoint of the exterior derivative.

\term{Hodge Theory}
The study of the cohomology of compact Kähler manifolds. The Hodge decomposition $H^k(X, \mathbb{C}) = \bigoplus_{p+q=k} H^{p,q}(X)$ relates the topology of the manifold to its complex structure.

\term{Hodge-Laplacian}
The operator $\Delta = d\delta + \delta d$ acting on differential forms. The kernel of this operator consists of the harmonic forms.

\term{Hofstadter's Butterfly}
A fractal pattern representing the energy spectrum of an electron in a 2D lattice subjected to a magnetic field.

\term{Hölder Space}
A space of functions whose differences are controlled by a power of the distance, i.e.\ functions satisfying $|f(x) - f(y)| \leq C |x - y|^\alpha$ for some $0 < \alpha \leq 1$.

\term{Hölder's Inequality}
The inequality $\int |fg| \leq \|f\|_p \|g\|_q$ for $1/p + 1/q = 1$, $p, q \geq 1$. It generalises the Cauchy--Schwarz inequality, which is the case $p = q = 2$.

\term{Holomorphic Function}
A complex-valued function of a complex variable that is complex-differentiable in a neighborhood of every point in its domain.

\term{Holonomy}
The group of linear transformations obtained by parallel transport of tangent vectors around closed loops on a manifold. It encodes curvature.

\term{Homeomorphism}
A continuous bijection $f: X \to Y$ whose inverse $f^{-1}$ is also continuous. Two spaces are homeomorphic if there exists a homeomorphism between them, meaning they are topologically identical.

\term{Homogeneous Coordinate}
A system of coordinates used in projective geometry to represent a point in $\mathbb{P}^n$ as a vector in $\mathbb{R}^{n+1}$, where two vectors represent the same point if they are scalar multiples.

\term{Homogeneous Polynomial}
A polynomial where every term has the same total degree. For example, $x^2 + xy + y^2$ is homogeneous of degree 2.

\term{Homogeneous Space}
A space with a transitive action of a Lie group, so that any point can be moved to any other by a group element. Examples include spheres and Grassmannians.

\term{Homological Algebra}
The branch of algebra studying chain complexes, exact sequences, and derived functors, unifying algebraic topology and homological invariants.

\term{Homology}
A tool in algebraic topology used to count the "holes" of different dimensions in a space. For a CW complex, the $n$-th homology group $H_n(X)$ describes the $n$-dimensional cycles that are not boundaries.

\term{Homomorphism}
A map between two algebraic structures (like groups, rings, or vector spaces) that preserves the operations. For groups $G$ and $H$, $f(ab) = f(a)f(b)$.

\term{Homotopy}
A continuous deformation between two continuous maps $f, g: X \to Y$. Two maps are homotopic if one can be continuously transformed into the other.

\term{Homotopy Category}
The category whose objects are topological spaces and whose morphisms are homotopy classes of continuous maps.

\term{Homotopy Group}
The group $\pi_n(X)$ of based homotopy classes of maps $S^n \to X$. The fundamental group $\pi_1$ is the case $n = 1$.

\term{Hopf Algebra}
A bialgebra equipped with an antipode, generalizing the group algebra of a group. They are central to the study of quantum groups and tensor categories.

\term{Hopf Fibration}
A fiber bundle $S^3 \to S^2$ where the fibers are circles $S^1$. It provides a way to view the 3-sphere as a collection of circles over a 2-sphere.

\term{Hopf Link}
The simplest non-trivial link, consisting of two circles interlocked once.

\term{Hopf-Rinow Theorem}
A theorem stating that for a Riemannian manifold, geodesic completeness is equivalent to metric completeness.

\term{Hopfian Group}
A group $G$ that is not isomorphic to any of its proper quotients. Every finite group is Hopfian.

\term{Horned Sphere}
A topological space that is homeomorphic to a 3-sphere but whose boundary is not a smooth 2-sphere, serving as a counterexample to several intuitive conjectures in topology.

\term{H-space}
A topological space $X$ equipped with a continuous binary operation that is homotopic to a group operation. Examples include loop spaces $\Omega Y$.

\term{Hurewicz Theorem}
A theorem relating the first non-vanishing homotopy group of a space to its homology group: if $\pi_i(X) = 0$ for $i < n$ with $n \geq 2$, then $\pi_n(X) \cong H_n(X)$.

\term{Hurwitz Zeta Function}
The function $\zeta(s, a) = \sum_{n=0}^\infty (n + a)^{-s}$, which generalises the Riemann zeta function $\zeta(s) = \zeta(s, 1)$.

\term{Hyperbolic Cosine (cosh)}
The hyperbolic function defined as $\cosh x = \frac{e^x + e^{-x}}{2}$. It satisfies $\cosh^2 x - \sinh^2 x = 1$.

\term{Hyperbolic Equation}
A second-order PDE whose characteristic equation has two distinct real characteristic directions, such as the wave equation. Hyperbolic problems admit well-posed initial value problems.

\term{Hyperbolic Function}
The functions $\sinh$, $\cosh$, and $\tanh$, defined in terms of exponentials and analogous to the trigonometric functions. They are the functions of a unit hyperbola rather than a circle.

\term{Hyperbolic Geometry}
A non-Euclidean geometry where the parallel postulate is replaced: through a point not on a line, there are at least two distinct lines parallel to the given line.

\term{Hyperbolic Group}
A finitely generated group whose Cayley graph is a hyperbolic metric space in the sense of Gromov.

\term{Hyperbolic Paraboloid}
A doubly ruled surface that looks like a saddle. It is defined by the equation $z = x^2/a^2 - y^2/b^2$.

\term{Hyperbolic Sine (sinh)}
The hyperbolic function defined as $\sinh x = \frac{e^x - e^{-x}}{2}$. It is the analogue of the sine function for hyperbolic geometry.

\term{Hyperbolic Space}
The standard model for hyperbolic geometry, such as the Poincaré disk or the upper half-plane. It has constant sectional curvature $-1$.

\term{Hyperbolic Tangent (tanh)}
The hyperbolic function $\tanh x = \frac{\sinh x}{\cosh x}$, used extensively as an activation function in neural networks.

\term{Hypercube}
The $n$-dimensional generalisation of a square and a cube, defined as the product of $n$ copies of the unit interval.

\term{Hyperelliptic Curve}
An algebraic curve of genus at least 2 that admits a degree-two map to the projective line. They generalise elliptic curves.

\term{Hypergeometric Function}
The function ${}_2F_1(a, b; c; z) = \sum_{n=0}^\infty \frac{(a)_n (b)_n}{(c)_n n!} z^n$, where $(x)_n$ is the Pochhammer symbol. It specialises to many elementary functions and orthogonal polynomials.

\term{Hypergraph}
A generalisation of a graph in which edges, called hyperedges, may connect any number of vertices.

\term{Hyperplane}
An $(n-1)$-dimensional subspace of an $n$-dimensional vector space. In $\mathbb{R}^3$, a hyperplane is a 2D plane.

\term{Hyperplane Arrangement}
A finite collection of hyperplanes in a vector space. The study of the topology of the complement of such an arrangement is central to the theory of braid groups.

\term{Hyperreal Number}
An extension of the real numbers that includes infinitesimals and infinite numbers, used in non-standard analysis.

\term{Hypersphere}
The set of all points in $\mathbb{R}^n$ at a constant distance from a fixed center.

\term{Hypersurface}
A manifold of dimension $n-1$ embedded in a manifold of dimension $n$.

\term{Hypervolume}
The $n$-dimensional volume of a region in $\mathbb{R}^n$ for $n > 3$.
